\section{Fixed-Form Scores from Shared Coordinates}\label{sec:pathdesign}
A form can reveal useful coordinates even when it fails to complete a tree path. For a frozen depth-$b$ tree, take the coordinate set on each root-to-leaf path, extend it to size $b$ with fixed unused coordinates, and retain the $L$ distinct sets as a tractable menu. Let $g_{T,y}(u)$ be a table on \emph{all} answers $u$ to form $T$, and let $\widehat P_T(u,y)$ be calibrated marginal frequencies. For binary features solve
\begin{align}
 \min_{q,g}\quad&\sum_{T,u,y}\widehat P_T(u,y)
                       \frac{g_{T,y}(u)^2}{q_T},\label{eq:fullprogram}\\
 \text{subject to}\quad&\sum_{T\supseteq S}\widehat g_{T,y}(S)
                   =y\widehat d(S)\quad\text{for every }S,y,\nonumber\\
 &\sum_Tq_T=1,\qquad q_T\ge\beta/L.\nonumber
\end{align}
Here $\widehat g_{T,y}(S)$ is the Fourier coefficient of the table on its own form; it is a linear transform of $g$. We use $\beta=0.1$ for fitted designs in the earlier experiments and $\beta=0.05$ in Section~\ref{sec:experiments}. These constraints are exactly equivalent to $\sum_Tg_{T,y}(x_T)=yd(x)$. Returning $g_{T,Y}(X_T)/q_T$ is therefore globally unbiased even when the reference frequencies are inaccurate; the reference affects only the variance.

There are at most $2L2^b$ score entries and coefficient constraints, and the program avoids enumerating the $2^m$ full profiles. Pairwise shared-coordinate corrections span the complete unbiased score space on the supplied menu over the full binary cube (Proposition~\ref{prop:sharing}). A correction adds $f(X_{T\cap T'},Y)/q_T$ on form $T$ and subtracts it divided by $q_{T'}$ on form $T'$. The optimization therefore includes every such control variate. In the language of coarsening at random, \eqref{eq:fullprogram} chooses the optimal augmentation of the inverse-probability-weighted score under the reference law, jointly with the form probabilities. For ordinal data, the implementation uses tensor contrasts on categories defined by the frozen tree's thresholds.

For an existing randomized archive, keep its known $q$ fixed and optimize only the tables. Form-specific marginal frequencies can be estimated from incomplete training records; they need not be assembled into a joint reference distribution. Identification then requires coverage by the actual form menu, not only the degree bound for the menu of all size-$b$ forms.

A \emph{completion-only} score returns an inverse-probability weighted target on a matching padded path and zero otherwise; the experiments call its uniform version path guessing. The \emph{prefix} comparator uses the answers along the tree up to its first unavailable query, discarding other coordinates on the form. With restricted pilot records, its leaf--label reference can be fitted by EM from the observed form, stopping node and label (the supplementary report).

\paragraph{Proper scoring rules.}
For Brier improvement, with $z=(Y+1)/2$,
\[
 (p_0-z)^2-(p_1-z)^2=d(1-p_0-p_1)+Yd.
\]
In the form program, replace $Yd$ by this known target and impose separate coefficient identities for each label. The same unbiasedness guarantee holds. For ensembles of degree $r$, Brier loss can require degree $2r$; log loss can introduce still higher interactions. Appendix~\ref{app:metrics} gives the identification conditions and explains why AUC requires a separate pairwise design.

\section{Choosing a Collection Design}\label{sec:guidance}
The resource constraint and available records determine the admissible comparison. A historical sample can estimate the target directly when it represents the target population; an external reference can calibrate scores without that transport assumption. Table~\ref{tab:decision} summarizes defaults.

\begin{table}[h]
\centering
\caption{Default choices. Training records are common historical inputs; new calibration records count toward the chosen budget.}
\label{tab:decision}
\small
\begin{tabular}{p{.45\textwidth}p{.47\textwidth}}
\toprule
Available information / constraint & Default audit \\
\midrule
Target-representative complete records or recorded predictions and label & Use their paired differences before collecting more records. \\
Hard cap; branching computes the contrast & Average $Yd$ from branching. \\
Hard cap; bias floor exceeds target half-width & Increase the cap, narrow a justified population class, or change the estimand. \\
Existing randomized forms; contrast covered & Keep the recorded allocation; fit scores from separate incomplete records and use the survey variance design. \\
Fixed forms required; feasible contrast; external reference & Fit pointwise-unbiased form scores; include any certified approximation bias. \\
Fixed forms; no usable calibration & Prespecify a positive design and unbiased score. Charge any pilot and test its cost sensitivity. \\
Occasional complete records allowed & Compare full-record averaging and stratified residual correction under the actual cost model. \\
Labels dominate cost; cap can be relaxed & Compare complete records or residual correction; saving a few features may save little total cost. \\
\bottomrule
\end{tabular}
\end{table}

For frozen imputation bias $\mathfrak b_I$ and variances $\sigma_I^2,\sigma_U^2$, imputation has smaller MSE exactly when $n\mathfrak b_I^2<\sigma_U^2-\sigma_I^2$. A fitted proxy can therefore win at small $n$ and lose as the audit grows with calibration held fixed. For inference about $\Delta$, use a certified bias bound or independent validation. A prospective bridge design should randomize collection mode and permitted item positions, and retain common anchors at matched positions, so that measurement and burden disparities can be detected \citep{schuman1981}.
